\section{Introduction}
\label{sec:intro}

Agentic systems for the sciences act on what their tools return. A copilot that answers a question about a protein, a planner that chooses a perturbation, or a closed-loop system that schedules an assay pulls records from databases such as UniProt \citep{uniprot2025} or PubChem \citep{kim2025pubchem} and reasons from them \citep{huang2025biomni,roohani2024biodiscovery,fa2026bioagent,bran2024chemcrow}. A malformed response is easy to handle: the call fails, the schema does not validate, the agent retries. A well-formed record with a wrong value is not. Nothing in the transport layer signals it, and the only defence is the agent's own domain knowledge. Does a scientific agent notice when a returned record is impossible? If it notices, does that change what it answers? And when it does correct a value, what is the correction made of?

Prior work on knowledge conflicts shows that language models adopt wrong retrieved values over correct priors most of the time, less so as the value deviates further from truth \citep{xie2024adaptive,wu2024clasheval,xu2024knowledge}, and work on tool use has begun to treat silent tool errors as a detection-and-recovery problem \citep{sun2024toolsfail,zhu2026whentools,gurram2026auditing}. Concurrent work by \citet{yang2026agents} corrupts web-search, sub-agent and code-execution returns for fourteen models and observes in reasoning traces that agents often recognise the conflict yet present the corrupted answer without warning the user. These studies infer whether the agent noticed from free reasoning text, or do not measure it at all, and none varies what the record offers the agent to check against. We make the detection channel explicit and gradable, and we manipulate the record.

We build a fully reproducible probe (Figure~\ref{fig:design}) in two domains: 300 reviewed human proteins from UniProt and 300 small molecules from PubChem; two mock lookup tools per domain; a corrupting layer that alters one field of one record per task and nothing else; and 300 tasks whose gold and corrupted-implied answers are computed by code and verified to differ. Ten models from five vendors (Anthropic, Google, OpenAI, DeepSeek, Moonshot) run every task under clean, impossible and plausible corruption. Three further manipulations vary what the agent can check: a severity sweep scales a field from $\times0.9$ to $\times0.005$ with its paired field either kept consistent or left unchanged; a direction split applies plausible corruptions as both shrinks and inflations; and a record-format switch replaces the seven-field summary with the full database entry, so that the corrupted value becomes derivable from the rest of the record. Every run is graded by exact or token matching, every trace is logged, and the whole study cost about \$235.

The results separate three behaviours that are usually conflated. First, \emph{detection without correction}: impossible values are flagged in 75--100\% of runs and still reach the final answer in 32--93\% of them, in both domains and for every vendor; the agent writes down why the value cannot be right and then reports it. Second, \emph{consistency, not plausibility}: what triggers detection is a contradiction inside the record or recall of the specific entity, not a general sense of what a protein or a molecule can be. An internally consistent record describing a five-residue, 500\,Da ``protein'' is answered from at 80--100\% by every model, while the same value paired with an unchanged mass is often repaired by recomputing one field from the other. Third, \emph{correction needs a route}: whether a flagged value is corrected is predicted field by field by whether the record contains something the value can be derived from. Give the agent the full UniProt entry and organism errors go from 100\% propagated to 0\%; give it the PubChem formula and 0 of 75 impossible numeric values survive, while fabricated names, which nothing in the record derives, still propagate in 64--80\%. Chemistry adds a refinement: because formulas are well memorised, models often detect a 20\% error by recalling the formula and recomputing---Opus~5 flags 96\% of plausible corruptions---and then relay the value in 98\% of those runs, while overriding impossible values in 54\%. Detection is not the bottleneck in either case; the model overrides the tool when the check is decisive and defers when it is not.

Our contributions are:
\begin{itemize}
\itemsep0pt
\item a deterministic, cheap probe of tool-value propagation with real records in two domains, byte-identical mock tools, code-verified gold answers, two record formats and no LLM judge (Section~\ref{sec:design});
\item evidence across ten models and five vendors that an explicit warning channel is used but does not move the answer, that plausible corruptions are absorbed regardless of direction, and that the model ordering transfers across domains, with task-clustered statistics (Sections~\ref{sec:dissociation}--\ref{sec:fields});
\item three manipulations---the consistency-controlled sweep, a cue classifier over 2{,}700 flagged concerns, and the record-format switch---that locate the correction that does occur in within-record derivation and specific recall rather than general plausibility, and show that the decline in adoption with implausibility reported by ClashEval vanishes when the record stays internally consistent (Sections~\ref{sec:cues}--\ref{sec:chem});
\item a measurement across vendors of what a one-sentence verification instruction buys, and what it does not (Section~\ref{sec:verify}).
\end{itemize}
